% NIH R01 — Data Management and Sharing Plan
% Required for all NIH applications since Jan 2023. Follow the six-element
% NIH DMS Plan structure.
\documentclass[11pt,letterpaper]{article}
\usepackage{nih-r01-support}
\begin{document}

\section{DATA MANAGEMENT AND SHARING PLAN}

\paragraph{Data type}
% TODO: Describe data types generated (e.g., synthetic / simulated, summary
%   statistics, derived results) and whether new individual-level human data
%   will be generated. State which data will be preserved and shared.

\paragraph{Related tools, software, and code}
% TODO: List the software / pipelines that will be released. Indicate
%   licenses (MIT/Apache 2.0). Note any specialized tools needed to access
%   the shared data.

\paragraph{Standards}
% TODO: Data and metadata standards (file formats, ontologies, GWAS Catalog
%   conventions, FAIR principles).

\paragraph{Data preservation, access, and associated timelines}
% TODO: Repositories (GitHub, Zenodo, dbGaP, GWAS Catalog, AnVIL, etc.) and
%   timeline relative to associated publications / end of award.

\paragraph{Access, distribution, or reuse considerations}
% TODO: Data use agreements (UK Biobank, dbGaP, All of Us), redistribution
%   restrictions, consent / privacy considerations, license under which
%   shared data are released.

\paragraph{Oversight of data management and sharing}
% TODO: Who oversees the DMS Plan (typically the PI), and for how long
%   after the award ends.

\end{document}
